\documentclass[10pt,a4paper]{amsart}
\usepackage[margin=19mm]{geometry}
\usepackage{amsmath,amssymb,mathtools}
\usepackage[scr=rsfs,cal=euler]{mathalfa}
\usepackage{xfrac}
\usepackage[unicode,hidelinks,bookmarks=false]{hyperref}
\newcommand{\ie}{i.e.\ }
\newcommand{\cf}{cf.\ }
\newcommand{\resp}{respectively\ }
\newcommand{\Subsection}{\subsection}
\providecommand{\nobreakdash}{\nobreak\hbox{-}}
\setlength{\emergencystretch}{2em}
\multlinegap0pt
\usepackage{bm}
\usepackage{bbm}
\usepackage{dsfont}
\usepackage{xspace}
\usepackage{cancel}
\usepackage{mathtools}
\usepackage{amsthm}
\makeatletter
\@ifundefined{theorem}{\newtheorem{theorem}{Theorem}}{}
\@ifundefined{lemma}{\newtheorem{lemma}[theorem]{Lemma}}{}
\@ifundefined{remark}{\newtheorem{remark}[theorem]{Remark}}{}
\makeatother
\makeatletter
\expandafter\ifx\csname claim\endcsname\relax
\newenvironment{claim}[1]{\par\noindent\underline{Claim:}\space#1}{}
\fi
\expandafter\ifx\csname claimproof\endcsname\relax
\newenvironment{claimproof}[1]{\par\noindent\underline{Proof:}\space#1}{\hfill $\square$}
\fi
\makeatother
\title[Low regularity well-posedness of nonlocal dispersive perturb]{Low regularity well-posedness of nonlocal dispersive perturbations of Burgers' equation}
\author{Luc Molinet, Didier Pilod, St\'ephane Vento}
\date{Source: arXiv:2506.17801v2 (CC BY 4.0)}
\begin{document}
\maketitle

\noindent\emph{Source and licence.} Low regularity well-posedness of nonlocal dispersive perturbations of Burgers' equation. Version arXiv:2506.17801v2 (CC BY 4.0), distributed under the Creative Commons Attribution 4.0 International licence, as recorded in the arXiv licence metadata for this exact version and shown on the abstract page https://arxiv.org/abs/2506.17801v2. Full rights notice: licenses/arxiv-2506.17801-CC-BY-4.0.txt.\par\smallskip

\noindent\textit{Editorial scope.} Counted unit: the Remark whose source label is \texttt{bound:a4tilde-resonant} in the article (source0.tex lines 1919--1934, source label \texttt{bound:a4tilde-resonant}), delivered together with the article's own complete original proof of it (source0.tex lines 1935--2095, one \texttt{proof} environment of 161 lines). No mathematics is written, repaired or re-derived: statement and proof are the article's own text line for line. Required context retained uncounted: def:a4tilde (lines 1007--1012); lem:b3tilde (lines 1838--1850); lem:b3tilde.1 (lines 1842--1844); lem:b3tilde.2 (lines 1846--1848); lem:diff-b3tilde (lines 1856--1861); lem:diff-b3tilde.0 (lines 1858--1860); a4tilde-B1 (lines 1884--1886); a4tilde-C2 (lines 1899--1901); a4tilde-C3 (lines 1907--1909); a4tilde-C4 (lines 1912--1914). Every definition, display and label of those lines is retained unchanged. Omitted from this package and not redistributed: 1--1006, 1013--1837, 1845--1845, 1849--1855, 1861--1883, 1887--1898, 1902--1906, 1910--1911, 1915--1918, 2096--2402. Nothing inside a retained excerpt is omitted or altered: every retained body is delivered exactly as the article's lines are written, so no omission receipt of any kind accompanies this package. Citations: the retained excerpts carry no citation command at all, so no bibliography entry of the article is transcribed and no reference is redistributed. Reference adaptation: none was needed and none was made; every reference inside the delivered unit resolves inside it, so no unresolved reference or citation is produced. Printed slips kept literal: no printed word, symbol, hypothesis or number of the counted statement or of its proof was altered. Layout and class: the article's own class is \texttt{amsart} with packages including babel, amssymb, amsmath, amsfonts, amsthm, enumerate, color, todonotes; the wrapper is the lane's pinned \texttt{amsart} editorial wrapper at 10pt on A4, which restates the theorem-like environments the retained text needs and adds the article's own macro definitions that the retained text needs (none), with extra wrapper packages (bm, bbm, dsfont, xspace, cancel, mathtools). The delivered numbering is the unit's own. Assets: no retained span contains a figure environment, a graphics-inclusion command, a PDF-page inclusion, a table or a TikZ picture; no image file is copied into this package, the package declares no asset dependency and no image file is redistributed with it. Licence: version arXiv:2506.17801v2 of this article is distributed under the Creative Commons Attribution 4.0 International licence, as recorded in the arXiv licence metadata for this exact version and shown on the abstract page https://arxiv.org/abs/2506.17801v2. A full rights notice is delivered at licenses/arxiv-2506.17801-CC-BY-4.0.txt. Characters: every retained excerpt is pure ASCII, so the delivered engine, XeLaTeX with its default Latin Modern text font, reports no missing character.
\begin{align}
\widetilde{a}_4(\xi_1, \xi_2, \xi_3) &= (\xi_1+\xi_3)\left[\widetilde{b_3}(\xi_1+\xi_3, \xi_2) - \widetilde{b_3}(\xi_1+\xi_3, -\xi_1)\right] \nonumber \\
&\quad + (\xi_2+\xi_3)\left[\widetilde{b_3}(\xi_2+\xi_3, \xi_1) - \widetilde{b_3}(\xi_2+\xi_3, -\xi_2)\right] \nonumber \\
&\quad - (\xi_1+\xi_2)\widetilde{b_3}(\xi_1, \xi_2) \nonumber \\
&=: (\widetilde{a}_{41} + \widetilde{a}_{42} + \widetilde{a}_{43})(\xi_1,\xi_2,\xi_3). \label{def:a4tilde}
\end{align}

\begin{lemma}\label{lem:b3tilde}
  Let $(\xi_1, \xi_2, \xi_3)\in \Gamma^3$ and $N_1,N_2,N_3$ be dyadic numbers such that $N_i\sim |\xi_i|$ for $i=1,2,3$. 
  \begin{enumerate}
    \item If $N_1\sim N_2$, one has
    \begin{equation}\label{lem:b3tilde.1}
    |\widetilde{b_3}(\xi_1, \xi_2)| \lesssim N_{max}^{2\sigma-\alpha}.
    \end{equation}
    \item If $N_1\ll N_2$ or $N_2\ll N_1$, one has
    \begin{equation}\label{lem:b3tilde.2}
    |\widetilde{b_3}(\xi_1, \xi_2)| \lesssim N_{min}^{-1} N_{max}^{2\sigma+1-\alpha}.
    \end{equation}
  \end{enumerate}
\end{lemma}

\begin{lemma}\label{lem:diff-b3tilde}
  With the notations of Lemma \ref{lem:b3tilde}, if $N_1\ll N_2$, then
  \begin{equation}\label{lem:diff-b3tilde.0}
    |\partial^{(0,1)}\widetilde{b}_3(\xi_1, \xi_2)| \lesssim N_1^{-1}N_2^{2\sigma-\alpha}.
  \end{equation}
\end{lemma}

      \begin{equation}\label{a4tilde-B1}
      |\widetilde{a}_4(\xi_1,\xi_2,\xi_3)| \lesssim N_1^{-1} N_{max}^{2\sigma+2-\alpha} .
    \end{equation}

       \begin{equation}\label{a4tilde-C2}
      |\widetilde{a}_4(\xi_1,\xi_2,\xi_3)| \lesssim N_1^{-1} N_{max}^{2\sigma+2-\alpha} .
    \end{equation}

   \begin{equation}\label{a4tilde-C3}
      \left|\widetilde{a}_4(\xi_1,\xi_2,\xi_3) + \frac{\xi_3\nu_\sigma(\xi_3)}{\xi_1\omega'(\xi_3)} \right| \lesssim N_1^{2\sigma}N_3^{1-\alpha} + N_1^{-1}N_3N_{max}^{2\sigma+1-\alpha} .
    \end{equation}

       \begin{equation}\label{a4tilde-C4}
      |\widetilde{a}_4(\xi_1,\xi_2,\xi_3)| \lesssim N_1^{2\sigma+1-\alpha} .
    \end{equation}

\begin{remark}
\begin{enumerate}
\item[$(i)$]
 For the region $C_2$ in the particular case $M_3:=|\xi_1+\xi_2|\ll N_1\sim N_2$,  we have the simpler bound
 $$
 |\widetilde{a}_4(\xi_1,\xi_2,\xi_3)| \lesssim N_1^{-1}N_2^{-1}M_3 N_{max}^{2\sigma+2-\alpha}.
 $$
 However, estimate \eqref{a4tilde-C2} is good enough for our purpose.
\item[$(ii)$]
In the resonant region $A$, we have the better bound 
\begin{equation} \label{bound:a4tilde-resonant}
|\widetilde{a}_{4}| \lesssim M_3N_{max}^{2\sigma-\alpha} .
\end{equation}
However, we will not use this bound in our energy estimate.
\end{enumerate}
\end{remark}

\begin{proof}
We will extensively use Lemma \ref{lem:b3tilde} to bound $\widetilde{a}_4$.
Recalling the definition of $\widetilde{a}_4$ in \eqref{def:a4tilde}, we also decompose 
\begin{equation}  \label{def:a4tilde.2}
\widetilde{a}_{41}=\widetilde{a}_{411}+\widetilde{a}_{412} \quad \text{and} \quad \widetilde{a}_{42}=\widetilde{a}_{421}+\widetilde{a}_{423}
\end{equation}
where 
\begin{align*}
\widetilde{a}_{411} & = (\xi_1+\xi_3) \widetilde{b_3}(\xi_1+\xi_3, \xi_2); \\
\widetilde{a}_{412}& =-(\xi_1+\xi_3) \widetilde{b_3}(\xi_1+\xi_3, -\xi_1); \\
\widetilde{a}_{421}& =(\xi_2+\xi_3)\widetilde{b_3}(\xi_2+\xi_3, \xi_1); \\
\widetilde{a}_{422}& = - (\xi_2+\xi_3)\widetilde{b_3}(\xi_2+\xi_3, -\xi_2) .
\end{align*}
We will denote by $M_1, M_2, M_3$ some dyadic numbers satisfying $M_1\sim|\xi_2+\xi_3|$, $M_2\sim |\xi_1+\xi_3|$ and, $M_3\sim |\xi_1+\xi_2|$.  
\vskip .3cm
\noindent
{\bf Region $A$}: $N_1\sim N_2\sim N_3\sim N_4 =: N$\\
Since $|\xi_1+\xi_3|\lesssim |\xi_1|\sim |\xi_2|$ Lemma \ref{lem:b3tilde} provides
\begin{align*}
|\widetilde{a}_{41}| &\le|(\xi_1+\xi_3)\widetilde{b}_3(\xi_1+\xi_3, \xi_2)| + |(\xi_1+\xi_3)\widetilde{b}_3(\xi_1+\xi_3, -\xi_1)|\\
&\lesssim   M_2(M_2^{-1}N^{2\sigma+1-\alpha} + M_2^{-1}N^{2\sigma+1-\alpha}) \\
&\lesssim N^{2\sigma+1-\alpha}.
\end{align*}
Similarly, we obtain the same bound for $|\widetilde{a}_{42}| + |\widetilde{a}_{43}|\lesssim N^{2\sigma+1-\alpha}$.


\vskip .3cm
\noindent
{\bf Region $B_1$}: $N_1\ll N_2\sim N_3\sim N_4=:N$\\
In this region, we have $M_1\sim M_2\sim M_3\sim N$. Therefore we get on one hand
$$
|\widetilde{a}_{411}| + |\widetilde{a}_{422}| + |\widetilde{a}_{43}| \lesssim N^{2\sigma+1-\alpha},
$$
and on the other hand
$$
|\widetilde{a}_{412}| + |\widetilde{a}_{421}| \lesssim N_1^{-1}N^{2\sigma+2-\alpha}.
$$
Then, \eqref{a4tilde-B1} follows immediately.

\vskip .3cm
\noindent
{\bf Region $B_2$}: $N_3\ll N_1\sim N_2\sim N_4 =: N$\\ 
Again, we have $M_1\sim M_2\sim M_3\sim N$ and every term $\widetilde{a}_{4ij}$, $i,j\in\{1,2\}$, as well as $\widetilde{a}_{43}$ is bounded by $N^{2\sigma+1-\alpha}$.

\vskip .3cm
\noindent
{\bf Region $C_1$}: $N_3\lesssim N_4 \ll N_1\sim N_2=: N$\\ 
In $C_1$ it holds that $M_3\lesssim N_4 \ll M_1\sim M_2\sim N$. First, we deduce from \eqref{lem:b3tilde.1} that
\begin{equation}\label{lem:a4tilde.1}
|\widetilde{a}_{43}| = |(\xi_1+\xi_2)\widetilde{b}_3(\xi_1,\xi_2)|  \lesssim M_3 N^{2\sigma-\alpha}.
\end{equation}
Using the symmetry in $(\xi_1,\xi_2)$ it remains to estimate 
\begin{align*}
\widetilde{a}_{411}+\widetilde{a}_{421} &= \xi_1[\widetilde{b}_3(\xi_1+\xi_3,\xi_2) - \widetilde{b}_3(\xi_2+\xi_3, \xi_1)] \\
&\quad +\xi_3 \widetilde{b}_3(\xi_1+\xi_3,\xi_2) - \xi_4 \widetilde{b}_3(\xi_2+\xi_3, \xi_1)\\
& := \widetilde{a}_{411}^{(1)} + \widetilde{a}_{411}^{(2)} + \widetilde{a}_{411}^{(3)}.
\end{align*}
From \eqref{lem:b3tilde.1} we get
\begin{equation}\label{lem:a4tilde.2}
|\widetilde{a}_{411}^{(2)}| + |\widetilde{a}_{411}^{(3)}| \lesssim (N_3+N_4)N^{2\sigma-\alpha} \lesssim N_4 N^{2\sigma-\alpha}.
\end{equation}
Then, we may rewrite $\widetilde{a}_{411}^{(1)}$ as
\begin{align*}
  \widetilde{a}_{411}^{(1)} &= \xi_1\left( \frac{\widetilde{a}_3(\xi_1+\xi_3,\xi_2)}{\Omega_3(\xi_1+\xi_3, \xi_2)} - \frac{\widetilde{a}_3(\xi_2+\xi_3, \xi_1)}{\Omega_3(\xi_2+\xi_3, \xi_1)} \right)\\
  &= \xi_1\Big( \widetilde{a}_3(\xi_1+\xi_3,\xi_2)\left(\frac{1}{\Omega_3(\xi_1+\xi_3, \xi_2)} - \frac{1}{\Omega_3(\xi_2+\xi_3, \xi_1)}\right) \\
  &\quad + \frac{\widetilde{a}_3(\xi_1+\xi_3,\xi_2) -\widetilde{a}_3(\xi_2+\xi_3, \xi_1)}{\Omega_3(\xi_2+\xi_3, \xi_1)}  \Big)\\
  &:= \widetilde{a}_{411}^{(11)} + \widetilde{a}_{411}^{(12)}.
\end{align*}
Using twice Taylor formula we deduce
\begin{align*}
 |\Omega_3(\xi_2+\xi_3, \xi_1) - \Omega_3(\xi_1+\xi_3, \xi_2)| &= |\omega_{\alpha+1}(\xi_2+\xi_3)-\omega_{\alpha+1}(\xi_2) + \omega_{\alpha+1}(\xi_1) - \omega_{\alpha+1}(\xi_1+\xi_3)|\\
 &\lesssim |\xi_3\omega_{\alpha+1}'(\xi_2) - \xi_3\omega_{\alpha+1}'(\xi_1)| + N_3^2N^{\alpha-1}\\
 &\lesssim N_3M_3N^{\alpha-1} + N_3^2N^{\alpha-1}\\
 &\lesssim N_4^2N^{\alpha-1}.
\end{align*}
It follows that
\begin{equation}\label{lem:a4tilde.3}
|\widetilde{a}_{411}^{(11)}| \lesssim N M_3N^{2\sigma} \frac{N_4^2N^{\alpha-1}}{(N_4N^\alpha)^2} \lesssim N_4N^{2\sigma-\alpha}.
\end{equation}
Similarly, recalling that $\widetilde{\nu}_\sigma(\xi)=\xi_|\xi|^{2\sigma}{\bf 1}_{|\xi|\ge 1}$ and $2\sigma+1>0$ we obtain
\begin{align*}
 |\widetilde{a}_3(\xi_1+\xi_3,\xi_2) -\widetilde{a}_3(\xi_2+\xi_3, \xi_1)| &\lesssim |\xi_4\widetilde{\nu}_\sigma'(\xi_1) - \xi_4\widetilde{\nu}_\sigma'(\xi_2)| + N_4^2N^{2\sigma-1}\\
 &\lesssim N_4M_3N^{2\sigma-1} + N_4^2N^{2\sigma-1}\\
 &\lesssim N_4^2N^{2\sigma-1},
\end{align*}
from which we conclude
\begin{equation}\label{lem:a4tilde.4}
|\widetilde{a}_{411}^{(12)}| \lesssim N\frac{N_4^2N^{2\sigma-1}}{N_4N^\alpha} \lesssim N_4N^{2\sigma-\alpha}.
\end{equation}
Gathering \eqref{lem:a4tilde.1}-\eqref{lem:a4tilde.2}-\eqref{lem:a4tilde.3}-\eqref{lem:a4tilde.4} we get the desired estimate for the region $C_1$.

\vskip .3cm
\noindent
{\bf Region $C_2$}: $N_1\lesssim N_2 \ll N_3\sim N_4=: N$\\ 
Thanks to \eqref{lem:b3tilde.2}, we directly obtain the rough bounds
$$
|\widetilde{a}_{411}| + |\widetilde{a}_{422}| \lesssim N_2^{-1}N^{2\sigma+2-\alpha} \lesssim N_1^{-1}N^{2\sigma+2-\alpha}
$$
and
$$
|\widetilde{a}_{412}| + |\widetilde{a}_{421}| \lesssim N_1^{-1}N^{2\sigma+2-\alpha}.
$$
In the same way, using \eqref{lem:b3tilde.1} or \eqref{lem:b3tilde.2} we always have
$$
|\widetilde{a}_{43}| \lesssim M_3 N_1^{-1}N_2^{-1} N_2^{2\sigma+2-\alpha}\lesssim N_1^{-1} N^{2\sigma+2-\alpha}.
$$

\vskip .3cm
\noindent
{\bf Region $C_3$}: $N_1\ll N_3 \ll N_2\sim N_4=: N$\\ 
We only show the most difficult estimate \eqref{a4tilde-C3}. Note that in $C_3$, we have $M_2\sim N_3$. According to Lemma \ref{lem:b3tilde} we get
$$
|\widetilde{a}_{411}| + |\widetilde{a}_{422}| \lesssim N^{2\sigma+1-\alpha} \lesssim N_1^{-1}N_3N^{2\sigma+1-\alpha}.
$$
For $\widetilde{a}_{412}$, we split
\begin{align}
\widetilde{a}_{412} &= -\xi_3\widetilde{b}_3(\xi_1+\xi_3,-\xi_1) + \mathcal{O}(N_3^{2\sigma+1-\alpha}) \notag \\
&= -\xi_3\frac{\nu_\sigma(\xi_1+\xi_3) - \nu_\sigma(\xi_1)}{\Omega_3(\xi_1+\xi_3, -\xi_1)} + \mathcal{O}(N_1^{2\sigma}N_3^{1-\alpha}). \label{lem:a4tilde.5}
\end{align}
Next, using a Taylor expansion of $\omega_{\alpha+1}$, we rewrite $$ \Omega_3(\xi_1+\xi_3, -\xi_1) = \xi_1\omega_{\alpha+1}'(\xi_3) + R(\xi_1,\xi_3)$$ with $|R(\xi_1,\xi_3)| \lesssim N_1^{\alpha+1}$. It follows that
\begin{equation}\label{lem:a4tilde.6}
\frac 1{\Omega_3(\xi_1+\xi_3, -\xi_1)} = \frac{1}{\xi_1\omega'_{\alpha+1}(\xi_3)} \frac{1}{1+\frac{R(\xi_1,\xi_3)}{\xi_1\omega'(\xi_3)}} = \frac{1}{\xi_1\omega'_{\alpha+1}(\xi_3)} + \mathcal{O}(N_1^{\alpha-1}N_3^{-2\alpha}) .
\end{equation}
Moreover, the mean value theorem provides
\begin{equation}\label{lem:a4tilde.7}
\widetilde{\nu}_\sigma(\xi_1+\xi_3) - \widetilde{\nu}_\sigma(\xi_1) = \widetilde{\nu}_\sigma(\xi_3) - \widetilde{\nu}_\sigma(\xi_1) + \mathcal{O}(N_1N_3^{2\sigma}) = \widetilde{\nu}_\sigma(\xi_3) + \mathcal{O}(N_1^{2\sigma+1}).
\end{equation}
Combining \eqref{lem:a4tilde.5}-\eqref{lem:a4tilde.6}-\eqref{lem:a4tilde.7} we infer
$$
\left|\widetilde{a}_{412} + \frac{\xi_3\widetilde{\nu}_\sigma(\xi_3)}{\xi_1\omega'(\xi_3)} \right| \lesssim N_1^{2\sigma}N_3^{1-\alpha}+N_3^{2+2\sigma-2\alpha}N_1^{\alpha-1} \lesssim N_1^{2\sigma}N_3^{1-\alpha},
$$
since $2\sigma+1-\alpha<0$.
To conclude the proof of \eqref{a4tilde-C3} we bound the term $\widetilde{a}_{421}+\widetilde{a}_{43}$ thanks to \eqref{lem:diff-b3tilde.0} by
\begin{align*}
  \widetilde{a}_{421}+\widetilde{a}_{43} &= \xi_2[\widetilde{b}_3(\xi_2+\xi_3,\xi_1) - \widetilde{b}_3(\xi_1,\xi_2)] + \mathcal{O}(N^{2\sigma+1-\alpha}) + \mathcal{O}(N_1^{-1}N_3N^{2\sigma+1-\alpha})\\
  &= \xi_2\xi_3\partial^{(0,1)}\widetilde{b}_3(\xi_1, \xi^\ast) + \mathcal{O}(N_1^{-1}N_3N^{2\sigma+1-\alpha}),\quad |\xi^\ast|\sim N\\
  &= \mathcal{O}(N_1^{-1}N_3N^{2\sigma+1-\alpha}).
\end{align*}

\vskip .3cm
\noindent
{\bf Region $C_4$}: $N_3\lesssim N_1 \ll N_2\sim N_4=: N$\\ 
From Lemma \ref{lem:b3tilde}, we easily check, by using $2\sigma+1-\alpha<0$, that
$$
|\widetilde{a}_{411}| + |\widetilde{a}_{422}| \lesssim M_2M_2^{-1}N^{2\sigma+1-\alpha} + N^{2\sigma+1-\alpha} \ll N_1^{2\sigma+1-\alpha}.
$$
Next, we claim that
$$
|\widetilde{a}_{412}| \lesssim N_1^{2\sigma+1-\alpha}.
$$
Indeed, if $N_3\ll N_1$, one has $M_2\sim N_1$ and $|\widetilde{b}_3(\xi_1+\xi_3, -\xi_1)| \lesssim N_1^{2\sigma-\alpha}$. In the case $N_3\sim N_1$, it holds $|\widetilde{b}_3(\xi_1+\xi_3, -\xi_1)| \lesssim M_2^{-1}N_1^{2\sigma+1-\alpha}$.\\
It remains to estimate
\begin{align*}
\widetilde{a}_{421} + \widetilde{a}_{43} &= 
\xi_2[\widetilde{b}_3(\xi_2+\xi_3,\xi_1) - \widetilde{b}_3(\xi_2, \xi_1)]  + [\xi_3\widetilde{b}_3(\xi_2+\xi_3,\xi_1) - \xi_1\widetilde{b}_3(\xi_2,\xi_1)]\\
&= \xi_2 \xi_3 \partial^{(1,0)} \widetilde{b}_3(\xi^\ast, \xi_1) + \mathcal{O}(N^{2\sigma+1-\alpha}), \quad |\xi^\ast|\sim N,\\
&= \mathcal{O}(N_3 N_1^{-1} N^{2\sigma+1-\alpha}) + \mathcal{O}(N_1^{2\sigma+1-\alpha})\\
&= \mathcal{O}(N_1^{2\sigma+1-\alpha}),
\end{align*}
where we used Lemma \ref{lem:diff-b3tilde}. This concludes the proof of \eqref{a4tilde-C4}.
\end{proof}

\end{document}
